% 第三章 硬件系统设计

\section{硬件系统设计}

\subsection{硬件总体架构}
\par 本系统硬件采用“采集节点 + 家庭网关 + 视觉节点 + 本地显示”的组合架构。采集节点以STM32为核心，负责环境传感器的驱动与数据采样；家庭网关以ESP32-S3为核心，负责多源数据汇聚、网络通信、显示与控制转发；视觉节点以OpenMV为核心，负责端侧人脸采集与识别；本地显示使用TFT屏实时展示关键环境指标与网络状态。该架构将“传感器采样”和“网络协议栈/交互显示”等不同实时性与复杂度的任务解耦，有利于降低单节点负载并提升系统可维护性。
\par 硬件连接关系上，STM32与ESP32-S3之间通过UART进行数据上报，串口波特率与数据格式在软件中统一约定。ESP32-S3与OpenMV之间通过独立UART通道进行命令与响应交互，支持文本响应与JPEG图像数据回传。ESP32-S3通过WiFi接入家庭路由器，向服务器端HTTP接口周期上报环境数据，同时通过MQTT通道接收控制指令并回传执行结果。TFT屏由ESP32-S3通过SPI接口驱动，实现本地可视化与状态提示。

\par 配图建议如下。图3-1建议为“硬件总体架构框图”，画法为模块框图：STM32采集节点（含I\(^2\)C传感器）、ESP32-S3网关（含TFT、WiFi、MQTT）、OpenMV视觉节点（含摄像头），并用连线标注UART1/UART2、I\(^2\)C、SPI、WiFi等接口。图3-2建议为“硬件连接关系逻辑图”，画法为更详细的连接图，标注具体引脚号（如UART1的TX/RX引脚、I\(^2\)C的SDA/SCL引脚、SPI的MOSI/SCK/CS引脚）与信号流向。
\begin{figure}[H]
\centering
\fbox{\parbox{0.92\linewidth}{\centering 图3-1占位：硬件总体架构框图。\\
要求展示STM32/ESP32-S3/OpenMV/TFT四大模块及其内部组成，标注接口类型（UART/I2C/SPI/WiFi）。}}
\caption{硬件总体架构框图（建议配图）}
\label{fig:hardware_architecture}
\end{figure}
\begin{figure}[H]
\centering
\fbox{\parbox{0.92\linewidth}{\centering 图3-2占位：硬件连接关系逻辑图。\\
要求标注具体引脚号（如GPIO17/18用于UART2，GPIO15/16用于OpenMV UART，PB6/PB7用于I2C等）与信号流向。}}
\caption{硬件连接关系逻辑图（建议配图）}
\label{fig:hardware_connections}
\end{figure}

\subsection{STM32传感器采集模块设计}
\par 采集模块以STM32为主控，采用I\(^2\)C总线连接AHT20温湿度传感器、BMP280气压传感器与SGP40 VOC传感器，实现对家庭环境的多维度感知。AHT20用于获取温度与湿度，BMP280用于获取气压并在必要时提供温度参考，SGP40在温湿度补偿输入的基础上输出VOC指数，适合用于空气质量趋势判断。传感器的供电与I\(^2\)C上拉需满足器件数据手册要求，原则上统一使用3.3V逻辑电平并保证供电去耦。
\par 软件设计采用“周期采样 + 结构化上报”的模式。在初始化阶段完成I\(^2\)C总线配置与各传感器的上电检测与初始化，并针对BMP280配置采样倍率与滤波参数以平衡噪声与响应速度。在运行阶段以固定周期读取温湿度、气压与VOC等数据，并对异常值（如读取错误或无效值）进行标记与容错处理。为降低耦合度，采集侧仅负责“产生可信数据”，不承担网络上传等复杂任务。
\par 数据通过USART2以串口文本形式上报至ESP32-S3。上报格式采用键值对组合并以换行符结束，便于网关侧按行解析与容错处理。典型格式包含温湿度字段（例如T与H）、气压字段（例如P）以及VOC字段（例如VOC），字段之间以逗号分隔并在行末结束，网关侧可根据字段存在性进行缺失值处理与融合。

\par 配图建议如下。图3-3建议为“STM32采集模块电路连接示意图”，画法为原理图风格：STM32主控、AHT20/BMP280/SGP40三个I\(^2\)C传感器、I\(^2\)C总线（SDA/SCL + 上拉电阻）、USART2接口（TX/RX），并标注I\(^2\)C地址与供电电压。图3-4建议为“STM32数据上报格式示例图”，画法为文本格式展示：例如“T:25.3,H:60.5,P:1013.2,VOC:125,”，并在下方标注字段含义与单位。
\begin{figure}[H]
\centering
\fbox{\parbox{0.92\linewidth}{\centering 图3-3占位：STM32采集模块电路连接示意图。\\
要求展示STM32、AHT20/BMP280/SGP40传感器、I2C总线（含上拉电阻）、USART2接口，标注I2C地址与供电。}}
\caption{STM32采集模块电路连接示意图（建议配图）}
\label{fig:stm32_circuit}
\end{figure}
\begin{figure}[H]
\centering
\fbox{\parbox{0.92\linewidth}{\centering 图3-4占位：STM32数据上报格式示例图。\\
要求展示典型数据行格式（如T:25.3,H:60.5,P:1013.2,VOC:125,），并在下方标注字段含义与单位。}}
\caption{STM32数据上报格式示例图（建议配图）}
\label{fig:stm32_data_format}
\end{figure}

\subsection{ESP32主控模块设计}
\par ESP32-S3在系统中承担“家庭网关”角色，核心任务是完成多源数据汇聚与网络侧双向通信。其一，ESP32-S3通过UART接收来自STM32的结构化文本数据，并解析更新系统内的环境数据结构。其二，ESP32-S3负责采集本地模拟量气体传感器（MQ系列）数据，以补充气体相关指标。其三，ESP32-S3通过WiFi连接家庭网络并以HTTP协议周期上报数据至服务器。其四，ESP32-S3集成MQTT客户端，实现控制指令订阅与响应发布，并将与OpenMV交互产生的文本/图像响应回传至平台侧。其五，ESP32-S3驱动TFT屏进行本地显示，以提升系统可用性与可观测性。
\par 软件架构采用模块化封装与分时调度策略。系统将UART接收、OpenMV通信、MQTT通信、传感器更新、屏幕刷新与HTTP上传分别作为独立任务，以非阻塞方式在主循环中轮询执行，并为不同任务设置独立的时间片间隔。例如，UART与MQTT处理以事件驱动为主，屏幕刷新与HTTP上传采用定时策略，从而避免网络请求或显示刷新对采样链路造成阻塞。该设计能够在资源受限条件下提供较稳定的端到端体验。
\par 显示界面设计以“关键信息优先”为原则，优先展示温度、湿度、气压、VOC与气体传感器读数等核心指标，并提供WiFi连接状态提示。通过固定刷新周期读取数据结构体并渲染到屏幕，可以减少界面闪烁并降低显示任务对系统实时性的影响。

\par 配图建议如下。图3-5建议为“ESP32网关功能模块框图”，画法为功能模块图：WiFi管理、MQTT通信、UART管理（接收STM32/发送OpenMV）、TFT显示、传感器采样（MQ气体）、数据结构管理，并用连线标注模块间数据传递关系。图3-6建议为“ESP32任务调度时序图”，画法为时序图：展示UART接收、传感器更新、显示刷新、HTTP上传、MQTT处理等任务的时间片分配与执行顺序。
\begin{figure}[H]
\centering
\fbox{\parbox{0.92\linewidth}{\centering 图3-5占位：ESP32网关功能模块框图。\\
要求展示WiFi/MQTT/UART/TFT/传感器采样/数据结构管理等模块，并用连线标注数据传递关系。}}
\caption{ESP32网关功能模块框图（建议配图）}
\label{fig:esp32_modules}
\end{figure}
\begin{figure}[H]
\centering
\fbox{\parbox{0.92\linewidth}{\centering 图3-6占位：ESP32任务调度时序图。\\
要求展示UART接收、传感器更新、显示刷新、HTTP上传、MQTT处理等任务的时间片分配与执行顺序。}}
\caption{ESP32任务调度时序图（建议配图）}
\label{fig:esp32_scheduling}
\end{figure}

\subsection{OpenMV人脸识别模块设计}
\par OpenMV视觉节点用于实现端侧人脸采集与识别能力。OpenMV具备摄像头输入与基础图像处理能力，适合在端侧进行轻量级视觉任务，从而减少敏感图像上云带来的隐私风险。系统在人脸检测阶段采用Haar级联分类器快速定位人脸区域，在识别阶段对人脸区域计算LBP特征并与本地样本库进行相似度匹配，通过阈值判定识别结果。该方法实现成本低、可在微型平台上运行，适合作为毕业设计的工程验证方案。
\par 功能实现以“命令驱动”的方式组织。OpenMV通过UART接收来自网关的控制命令并进入不同工作模式，包括人脸采集、识别、状态查询、图像获取、名单列举与删除等。采集模式下系统按固定数量保存人脸样本到本地存储目录，识别模式下系统持续输出识别结果文本消息，图像获取命令则触发OpenMV抓拍并以JPEG二进制流回传。为提升系统稳定性，OpenMV端对内存使用与异常情况进行处理，并在必要时回退到空闲模式。
\par 人脸数据库采用文件系统目录组织，以“姓名目录 + 样本文件”的方式保存采集到的人脸图像样本。系统在进入识别模式时遍历数据库并加载样本特征到内存，以降低运行时的重复计算开销。该设计便于在不改动协议的情况下扩展更多用户样本，同时便于实现“列出名单、删除用户”等管理功能。

\par 配图建议如下。图3-7建议为“OpenMV人脸识别流程示意图”，画法为流程图：图像采集→Haar检测→LBP特征提取→与数据库匹配→阈值判定→输出结果，并在关键步骤标注算法名称（Haar级联、LBP特征、相似度匹配）。图3-8建议为“OpenMV命令与响应交互示意图”，画法为状态机或流程图：展示空闲/采集/识别三种模式切换，以及各模式下的命令输入与响应输出。
\begin{figure}[H]
\centering
\fbox{\parbox{0.92\linewidth}{\centering 图3-7占位：OpenMV人脸识别流程示意图。\\
要求展示从图像采集到识别结果输出的完整流程，标注Haar检测、LBP特征提取、相似度匹配等关键步骤。}}
\caption{OpenMV人脸识别流程示意图（建议配图）}
\label{fig:openmv_flow}
\end{figure}
\begin{figure}[H]
\centering
\fbox{\parbox{0.92\linewidth}{\centering 图3-8占位：OpenMV命令与响应交互示意图。\\
要求展示空闲/采集/识别三种模式的状态转换，以及各模式下的命令输入与响应输出。}}
\caption{OpenMV命令与响应交互示意图（建议配图）}
\label{fig:openmv_commands}
\end{figure}

\subsection{硬件接口与交互设计}
\par 接口设计以可靠性与可调试性为原则。STM32与ESP32-S3之间的UART链路采用常用的8N1格式与115200波特率，能够满足环境数据的周期上报需求并便于串口调试。ESP32-S3与OpenMV之间使用独立UART通道以隔离数据上报与控制交互，避免不同业务数据混线导致解析复杂度上升。
\par 传感器与STM32之间采用I\(^2\)C总线连接，需保证总线时序满足器件要求，并在硬件上配置合适的上拉电阻与供电去耦。由于系统采用多器件共享I\(^2\)C总线，需要确保地址不冲突，并在软件初始化阶段对设备在线状态进行检测与容错处理。
\par 供电与电平方面，系统各模块原则上使用3.3V逻辑电平以保证互联兼容性。对于可能存在5V供电的外设，应通过电平转换或选用3.3V兼容模块避免直接连接导致的可靠性问题。为提升系统稳定性，建议在电源输入端与关键芯片附近布置去耦电容，并在电源设计上考虑瞬态电流、纹波与地线回流路径，以减少对模拟采样与I\(^2\)C通信的干扰。
